\documentclass[12pt,a4paper]{report}
\usepackage[utf8]{inputenc}
\usepackage[T1]{fontenc}
\usepackage[english,italian]{babel}
\usepackage{graphicx}
\usepackage{csquotes}
\usepackage{amsmath}
\usepackage{hyperref}
\usepackage{natbib}
\usepackage{booktabs}

% ************************ CODE ************************
\usepackage{listings} % Add this for code listings

% Configuring the listings package to handle Python code
\usepackage{color}
\definecolor{codegreen}{rgb}{0,0.6,0}
\definecolor{codegray}{rgb}{0.5,0.5,0.5}
\definecolor{codepurple}{rgb}{0.58,0,0.82}
\definecolor{backcolour}{rgb}{0.95,0.95,0.92}

\lstdefinestyle{mystyle}{
    backgroundcolor=\color{backcolour},   
    commentstyle=\color{codegreen},
    keywordstyle=\color{magenta},
    numberstyle=\tiny\color{codegray},
    stringstyle=\color{codepurple},
    basicstyle=\footnotesize\ttfamily,
    breakatwhitespace=false,         
    breaklines=true,                 
    captionpos=b,                    
    keepspaces=true,                 
    numbers=left,                    
    numbersep=5pt,                  
    showspaces=false,                
    showstringspaces=false,
    showtabs=false,                  
    tabsize=4
}

\lstset{style=mystyle}
% ************************ CODE ************************


% Define paths to images and other resources
\graphicspath{{images/}} % Assuming your logo is in the images folder

% Document metadata
\title{Your Thesis Title Here}
\author{Your Name Here}
\date{Date of Submission}

% PDF metadata and links configuration
\hypersetup{
    pdftitle={Your Thesis Title Here},
    pdfauthor={Your Name Here},
    pdfsubject={Master's Thesis in Artificial Intelligence and Cybersecurity},
    colorlinks=false, % Change to true if you prefer colored links
    linkcolor=blue,
    citecolor=blue,
    filecolor=magenta, 
    urlcolor=cyan
}

\makeatletter
% Customize the appearance of the front page
\newcommand{\frontpage}{
    \begin{titlepage}
        \centering
        \includegraphics[width=0.3\textwidth]{logo} % Replace 'logo' with 'logo.png' or 'logo.pdf'
        \vspace{1cm}
        
        \textbf{\Large{Kore University of Enna}}\\
        \vspace{2cm}
        \textbf{\Large{Master Degree in}}\\
        \vspace{0.2cm}
        \textbf{\Large{Artificial Intelligence and Cybersecurity}}\\
        \vspace{2cm}
        
        \textbf{\Huge{\@title}}
        \vspace{1.5cm}
        
        \begin{flushright}
            \textbf{Author:}\\
            \vspace{0.2cm}
            \textbf{\@author}
            \vspace{1cm}
        \end{flushright}

        
        \begin{flushleft}
            \textbf{Supervisor: } Nome del Relatore\\
            \vspace{0.2cm}
            \textbf{Co-Supervisor: } Nome del Correlatore\\
            \vspace{2cm}
        \end{flushleft}
        
        \textbf{\@date}
        \vfill
    \end{titlepage}
}
\makeatother

\begin{document}
\selectlanguage{english} % Change this to 'italian' if you want the base language of your document to be Italian

\frontpage

\chapter*{Abstract}
\addcontentsline{toc}{chapter}{Abstract}
This is the abstract of the thesis, where you should summarize the content of the thesis in about 200 words.

\chapter*{Sommario}
\addcontentsline{toc}{chapter}{Sommario}
Questo è il sommario della tesi, dove si riassume il contenuto della tesi in circa 200 parole.

\tableofcontents % Automatically generate a table of contents
\listoffigures   % Automatically generate a list of figures

\chapter{Introduction}
\label{ch:introduction}

This contains some introductory text.

\section{Background}
\label{sec:background}
Text for the background subsection may go here.

\begin{figure}[ht] % Positioning option [htbp] means here, top, bottom, page
\centering
\includegraphics[width=0.8\textwidth]{logo.png} % Replace 'example-image' with your actual image file name
\caption{This is an example of an image with a caption.}
\label{fig:example}
\end{figure}

An example of image can be found in Figure \ref{fig:example}.

\chapter{Literature Review}
\label{ch:literature_review}

Discuss relevant literature here.
Example of citation using te associated bib file \cite{rosenblatt1958perceptron}.

\chapter{Methodology}
\label{ch:methodology}
Describe your research methods here.

\section{Example Algorithm}
Below is the Python code for generating a Fibonacci sequence up to a specified number \( n \):

\begin{lstlisting}[language=Python, caption=Fibonacci Sequence in Python]
def fibonacci(n):
    a, b = 0, 1
    while a < n:
        print(a, end=' ')
        a, b = b, a+b
    print()

# Example of calling the function
fibonacci(100)
\end{lstlisting}

\chapter{Results}
\label{ch:results}

Present your findings here with tables similar to Table \ref{tab:results} and Figures similar to \ref{fig:example}.

\begin{table}[htbp]
\centering

\begin{tabular}{@{}lcc@{}} % lcc = left, center, center
\toprule
\textbf{Item} & \textbf{Quantity} & \textbf{Unit Price} \\
\midrule
Apples & 4 & \$1.50 \\
Oranges & 10 & \$2.00 \\
Bananas & 5 & \$1.75 \\
Grapes & 3 & \$2.50 \\
\bottomrule
\end{tabular}
\caption{Example of a Table - you may want to use \url{https://www.tablesgenerator.com/} for simplicity \textbf{with booktabs style!}}
\label{tab:results}
\end{table}

\chapter{Discussion}
\label{ch:discussion}
Discuss the implications of your results here.

\chapter{Conclusion}
\label{ch:conclusion}
Conclude your thesis and...

\section{Future Works}
\label{sec:future_works}

Discuss possible future research directions.


\appendix
\chapter{Additional Material}
Any additional supporting material can be included here.

\bibliographystyle{plain}
\bibliography{bibliography}


\end{document}
